\section{Introduction}\label{sec:intro}

\subsection*{Background}
\emph{Exactly divergence-free elements.} Scott and Vogelius \cite{ScottVogelius1985} showed that continuous piecewise-$P_k$ velocities with discontinuous piecewise-$P_{k-1}$ pressures satisfy $\div\Vh=\Pih$ for $k\ge4$ on meshes without singular vertices; the inf-sup constant was later made explicit in terms of the local mesh geometry \cite{GuzmanScott2019}. The discrete velocity is then divergence-free pointwise, not only weakly. This gives exact mass conservation. On fitted meshes with strongly imposed no-slip data it also gives pressure robustness: the velocity error does not depend on the pressure, which matters when pressure gradients are large or the viscosity is small \cite{JohnEtAl2017}. Split meshes give similar pairs of lower degree \cite{GuzmanNeilan2018}; see \cite{Neilan2020} for a review.

\emph{Curved walls.} Exact incompressibility has a cost at a curved wall. A divergence-free field that is smooth at a polygon vertex and vanishes on both adjacent edges has zero gradient there, and piecewise-polynomial divergence-free fields inherit such gradient constraints at the vertices of a polygonal wall \cite{ScottPrize,GjerdeScott2024}. If the wall is replaced by an inscribed polygon and no-slip is imposed strongly, the computed gradient can therefore be off by $O(1)$ at the polygon vertices, and the $H^1$ error of Scott--Vogelius elements then decays only like $h^{1/2}$ \cite{ScottPrize,GjerdeScott2024}, compared with $h^{3/2}$ for scalar problems \cite{BergerScottStrang1972,Scott1975}. How strongly this locking acts depends on the mesh at the wall, at least in our computations (Section~\ref{num:strong}). For a conforming piecewise-polynomial field the constraints at a wall vertex depend on its star: with two triangles they force a zero gradient, with three they leave a one-parameter family. On a mesh in which alternate polygon vertices lie in only two triangles we reproduce the rate $h^{1/2}$ with an $O(1)$ vertex-gradient error. On the meshes used in the rest of this paper, where every polygon vertex lies in three triangles, strong imposition shows no locking in the range computed ($H^1$ rates decreasing from $1.93$ to $1.61$ at $N=128$). We have not analysed strong imposition, and we do not claim that three triangles per wall vertex suffice in general; the observation suggests that the $h^{1/2}$ phenomenon needs wall vertices whose stars are singular or nearly singular.

One remedy is to fit the geometry, with isoparametric or Piola-mapped elements \cite{NeilanOtus2021,DurstNeilan2024}. Gjerde and Scott \cite{GjerdeScott2024} keep the polygon and impose the wall condition weakly by Nitsche's method \cite{Nitsche1971}, with penalty $\mu/h$ and no pressure term, and report errors orders of magnitude below those of strong imposition. Scott observed $H^1$ errors of order $h^{3/2}$ for this Scott--Vogelius--Nitsche method on a benchmark with constant pressure \cite{ScottPrize}. His zero-gradient prize \cite{ScottPrize} asks for a published proof that Scott--Vogelius--Nitsche converges, including a computational study of the influence of the Nitsche penalty. It states the expected error $\hG^{3/2}+\hO^k$ and adds: \emph{if this does not hold, why not?} We treat the question for general Stokes data.

\subsection*{Answer}
For the method as printed in \cite[(27)--(28)]{GjerdeScott2024} and bounded $\gamma$, the expected estimate holds if and only if the pressure is constant on the wall. The printed Nitsche form contains no traction term $-pn$. When $p$ varies along $\Gamma$, the penalty, not a traction, balances the pressure. The discrete velocity then leaks through the wall with normal velocity $(h/\mu)(p-\pbar)$, and the error is dominated by this leak (Theorems~A--C, Corollary~B$'$ and Theorem~B$''$). The mean-free correction $\CNS$ (cf.\ Frachon, Nilsson and Zahedi \cite{FrachonNilssonZahedi2024}), well posed for $\gamma\ge\gamma_2$, restores the expected rate (Theorem~D), and its $\hG^{3/2}$ term cannot be improved whenever the wall shear is not identically zero (Theorem~D$'$). On every mesh whose shortest boundary edge carries a triangle with apex angle at least $35^\circ$ and base angles at least $5^\circ$, the penalty threshold must grow with the ratio of bulk to boundary mesh size (Proposition~E).

\subsection*{Main results}
The notation is fixed in Section~\ref{sec:setting}. Throughout, $h=\hO$ is the bulk mesh size used in Scott's penalty $\mu/h$, $\hG$ is the length of the longest boundary edge, $\rho=h/\min_e|e|$, $\gamma=\mu\hG/h$ and $G=\norm{p-\pbar}_{L^2(\Gamma)}$. Constants do not depend on $h,\hG,\mu,\gamma$ or $\rho$.

\begin{theoremA}[Theorem~\ref{thm:A}]
The printed method $\GS$ converges, and its error in the mesh-dependent energy norm is
\[
\tnorm{\ut-u_h}\;\asymp\;(h/\mu)^{1/2}\,G\;+\;O\bigl((1+\gamma^{1/2})\hG^{3/2}+\eps\bigr),
\]
with matching upper and lower bounds.
\end{theoremA}

\begin{theoremC}[Theorem~\ref{thm:C}]
$\norm{\nabla(\ut-u_h)}\le C_p\,h/\mu+C\bigl((1+\gamma^{1/2})\hG^{3/2}+\eps\bigr)$.
The rate $1$ (rather than $1/2$) holds because $\div v=0$ turns the normal load into a tangential derivative (Lemma~\ref{lem:cancel}).
\end{theoremC}

\begin{theoremB}[Theorem~\ref{thm:B}]
In the regime where the pressure term dominates, the slip satisfies $\norm{\ut-u_h}_{L^2(\Gh)}\asymp (h/\mu)G$ and $\norm{\nabla(\ut-u_h)}\asymp h/\mu$. So Theorem~C is sharp.
\end{theoremB}

\begin{corollaryBp}[Corollary~\ref{cor:Bp}]
$\ut-u_h=(h/\mu)\,v_\varphi+o\bigl((h/\mu)^{1/2}\bigr)$ in the energy norm, where $v_\varphi$ is an explicit divergence-free field with $v_\varphi\approx-(p-\pbar)n$ on $\Gh$. In particular
\[
\frac{\norm{\ut-u_h}_{L^2(\Gh)}}{(h/\mu)G}\to1,
\qquad
\frac{\tnorm{\ut-u_h}}{(h/\mu)^{1/2}G}\to1.
\]
\end{corollaryBp}

\begin{theoremBpp}[Theorem~\ref{hc:thm}, Corollary~\ref{hc:cor}]
Let $U$ be the Stokes flow in $\Omega$ with Dirichlet data $-(p-\pbar)n$ on $\Gamma$ and $0$ on $\partial R$. For fixed $\gamma$ and bounded $\rho$, $(\mu/h)(\ut-u_h)\to U$ in $H^1$, so
\[
\frac{\norm{\nabla(\ut-u_h)}}{(h/\mu)G}\to K:=\frac{\norm{\nabla U}_{L^2(\Omega)}}{G}.
\]
The constant $K$ is not universal: it depends on the domain. For the pressure of our test P it is $K=2.7565157$ (series computation, inside a rigorous bracket $[2.0746,3.0526]$), and $K\to3/\sqrt7$ as the outer box grows. Convergence in the regime $\gamma\hG\to\infty$ is observed but not proved.
\end{theoremBpp}

\begin{theoremD}[Theorems~\ref{thm:D}, \ref{rs:unif} and~\ref{rs:Dfc}]
The naive consistent form, which adds $\ip{p_h}{v\cdot n}$, is singular; its $L^2_0$ variant is square but only gives $\div u_h\equiv c$, with $c\ne0$ in general (Proposition~\ref{prop:singular}). The mean-free correction $\CNS$, which adds $\ip{p_h-\pbG(p_h)}{v\cdot n}$ to the momentum row only, is well posed for $\gamma\ge\gamma_2$ and keeps $u_h$ exactly divergence-free; its velocity coincides with that of the zero-net-flux formulation of \cite{FrachonNilssonZahedi2024} (Remark~\ref{rem:fnz}). With the flux-corrected outer data $\tilde g_I$ of Definition~\ref{rs:def}, for general Stokes data, every fixed $k\ge4$ and every $\gamma\ge\gamma_2$,
\[
\norm{\ut-u_h}_{H^1(\Oh)}\le C\bigl(\hG^{3/2}+\hO^{k}\bigr),
\]
with $C$ independent of $\gamma$ and with no relation between $\hG$ and $\hO$. This is Scott's expected rate. With the plain interpolant $\gI$ the same holds under the proviso that the boundary edges are not too short relative to the bulk ($\hG\ge\hO^4$ for even $k$, $\hG\ge\hO^2$ for odd $k$); without it a single lower-order term, due to the net flux of the interpolated outer data, remains. On Clough--Tocher refinements the estimate of Theorem~\ref{thm:D} holds for every $k\ge2$ (Theorem~F).
\end{theoremD}

\begin{theoremDp}[Theorem~\ref{lb:thm}, Corollary~\ref{lb:cor}]
Suppose the wall shear $W=\dn u\cdot\tau$ is not identically zero. For $\CNS$, and for $\GS$ with any $G$,
\[
\norm{\nabla(\ut-u_h)}\ge c\,\norm{W}_{L^2(\Gamma)}\,\hG^{3/2}-C\hG^2 ,
\]
with $c$ independent of $\gamma,\mu,\rho$ and the pressure. So, whenever $\eps\le C\hG^{3/2}$ and $\gamma$ is bounded (or the data are flux-corrected and $\hO^k\le\hG^{3/2}$), the rate $\hG^{3/2}$ of Theorem~D (and of Theorem~A at $G=0$) is exact; it is the $\hG^{3/2}$ term, not the $\hO^k$ term, that is shown to be sharp. This is \emph{unconditional for $k\ge7$}. \emph{For $4\le k\le6$ it is conditional} on a local, scale-invariant condition (M3) on the boundary stars, which we have not proved; we verified it numerically at every boundary vertex of our meshes, for $N$ up to $1024$.
\end{theoremDp}

\begin{propositionE}[Proposition~\ref{prop:penalty}, Theorem~\ref{pb:thm}, Corollary~\ref{pb:rho}]
The condition $\mu\ge4\kappa C_I^2\rho$, $\rho=h/\min_{e\in\EG}|e|$, is sufficient for coercivity on every mesh. On every mesh, coercivity requires $\mu\ge\lambda_T\,h/|e|$ for each triangle $T$ with an edge $e$ on $\Gh$ (and not touching $\partial R$), with an explicit $\lambda_T$ depending only on the shape of $T$. In exact rational arithmetic, $\lambda_T>1$ whenever the apex angle is at least $35^\circ$ and $\lambda_T>9.3$ when it is at least $60^\circ$ (base angles at least $5^\circ$). So $\mu\gtrsim\rho$ is necessary on every mesh whose shortest boundary edge carries a triangle with apex angle at least $35^\circ$ and base angles at least $5^\circ$; for tall boundary triangles $\lambda_T<0$, and a necessity statement under (M0)--(M2) alone is open. Necessity is also certified on stars close to two reference stars.
\end{propositionE}

Theorem~A with $G=0$ gives $\tnorm{\ut-u_h}\le C((1+\gamma^{1/2})\hG^{3/2}+\eps)$. This recovers, for the printed method, the constant-pressure benchmark studied numerically in \cite{GjerdeScott2024,ScottPrize} and analysed in \cite{Cavalcante}, where the bound $\norm{u-u_h}_{H^1}\le C(\hO^k+\hG^{3/2})$ is proved for that benchmark, for fixed $k\ge4$ and also for $k\ge2$ on Clough--Tocher splits \cite[abstract and Cor.~5.3]{Cavalcante}. Our proof does not use \cite{Cavalcante}. Theorems~A--D$'$ hold for every fixed $k\ge4$ on meshes satisfying (M0)--(M2).

\subsection*{Further results}
\begin{theoremR}[Theorems~\ref{rs:unif}, \ref{rs:elower}, Corollaries~\ref{rs:rescue}, \ref{rs:rates}, Proposition~\ref{rs:cond}]
The $H^1$ error of the printed method is bounded uniformly in the penalty: $\norm{\nabla(\ut-u_h)}\le C_p\,\hG/\gamma+C(\hG^{3/2}+h^k+h^{\sigma_k}\hG^{-1/2})$ for every $\gamma\ge\gamma_0$. So a growing penalty $\mu\gtrsim h\hG^{-3/2}$ ($\mu\gtrsim h^{-1/2}$ when $h\asymp\hG$) restores the $H^1$ rate $\frac32$ of $\GS$ for general data. It cannot restore the energy norm: if $W\not\equiv0$ the energy error is at least $c\norm{W}\gamma^{1/2}\hG^{3/2}-C(\hG/\gamma)^{1/2}$ for large $\gamma$, so under a power-law penalty $\gamma\asymp\hG^{-\alpha}$ the energy rate of $\GS$ is $\min(\frac{1+\alpha}2,\frac{3-\alpha}2)\le1$ when $G>0$ and $W\not\equiv0$. The condition number of the velocity block grows by the factor $1+\gamma$, that of the pressure Schur complement by the factor $1+\mu/h$.
\end{theoremR}

\begin{theoremP}[Section~\ref{pd:sec}]
Let $p^\circ=\pt-\pbG(\pt)$. For $\CNS$ with flux-corrected data and $\gamma\in[\gamma_2,\gamma_{\max}]$, $\norm{p^\circ-p_h}_{L^2(\Oh)}\le C(\hG^{3/2}+\hO^k)$, including the constant, with $C$ depending on $\gamma_{\max}$ (Theorem~\ref{pd:thm:cnsp}). For $\GS$ the pressure is uniquely determined, but carries an explicit first-layer field $\xib$ produced by the symmetric Nitsche term acting on the leak; for fixed $\gamma$, bounded $\rho$ and $G>0$ the $L^2$ pressure error converges at the exact rate $\frac12$, and not at all pointwise in the first layer (Theorem~\ref{pd:thm:gsp}). The variational drag, lift and torque coincide with the integral of the Nitsche flux. For $\GS$ with fixed $\gamma$ and bounded $\rho$ their error is $-(h/\mu)\int_\Gamma(p-\pbar)(R_\psi-\bar R_\psi)+O(h^{3/2})$, with $R_\psi$ the pressure of a dual Stokes problem, so the rate is $1$ (Theorem~\ref{pd:thm:gsdrag}); for $\CNS$ with $\gamma\ge\max(1,\gamma_2)$ and $\gamma\hG\le1$ it is $O(\hG^2+\eps)$ (Theorem~\ref{pd:thm:cnsdrag}).
\end{theoremP}

\subsection*{Related work and what is new}
For Stokes, the consistent Nitsche form uses the full traction $\sigma(u,p)n$. It therefore contains the pressure term $\ip{p}{v\cdot n}$ \cite{FreundStenberg1995,Stenberg1995,Becker2002,BurmanHansbo2014}, and this term is recognised as essential for consistency \cite{BurmanHansboLarson2019Stokes}. The Gjerde--Scott form omits it. On the geometric side, boundary-value correction \cite{BrambleDupontThomee1972} has been combined with Nitsche's method for cut elements \cite{BurmanHansboLarson2018,BurmanHansboLarson2019Stokes} and for polygonal approximations of scalar problems \cite{DupontGuzmanScott2025}; the shifted boundary method \cite{MainScovazzi2018} is a related approach. None of these works uses exactly divergence-free spaces.
%% VERIFY: FreundStenberg1995 treats Stokes only per a secondary citation; Becker2002 content not read (notes/v5/lit_report.md).

With such spaces, the placement of the pressure term decides whether exact incompressibility survives. If the term also appears in the continuity equation, $\div u_h$ vanishes only away from an $O(h)$ boundary strip \cite{LiuNeilanOlshanskii2023}. If it appears only in the momentum equation, the constant pressure mode must be handled separately. Frachon, Nilsson and Zahedi \cite{FrachonNilssonZahedi2024} show that fixing the pressure mean by a Lagrange multiplier makes $\div u_h$ equal to that multiplier, and avoid this by testing the momentum equation only with velocities of zero net flux. Our $\CNS$ is this device in square form: its velocity coincides with theirs (Remark~\ref{rem:fnz}), and its wall-mean subtraction fixes the pressure constant automatically. The work closest to ours is that of Liu, Neilan and Otus \cite{LiuNeilanOtus2023}, for Scott--Vogelius elements on Clough--Tocher splits of a mesh inside a curved domain. They add a mean-free boundary multiplier that approximates the wall pressure up to a constant, together with a zero-flux constraint and a Taylor boundary correction. Their velocity is exactly divergence-free and converges at the optimal order $h^k$; without the constraint, they note, the method is ill posed \cite[Remark~3.2]{LiuNeilanOtus2023}. Related multiplier formulations for cut elements are given in \cite{BurmanHansboLarson2024}. Exactly divergence-free methods also require discrete Dirichlet data with zero net flux \cite{EickmannScottTscherpel2025}; our flux-corrected $\tilde g_I$ is a simple instance.

Our subject is the method as Gjerde and Scott printed it, and its minimal repair: no extra unknowns, no boundary correction, and the pressure itself on the wall. We do not claim the mean-free device as new. As far as we know, what is new is the analysis:
\begin{itemize}
\item the two-sided leak analysis of the printed method (Theorems~A--C, Corollary~B$'$), including the identification of its $H^1$ constant (Theorem~B$''$) and of its pressure and drag defects (Theorem~P);
\item the proof of Scott's conjectured rate $\hG^{3/2}+\hO^k$ for the mean-free correction on an inscribed polygon (Theorem~D), with the penalty-uniform $H^1$ bound, and the proof that its $\hG^{3/2}$ term is sharp (Theorem~D$'$);
\item the penalty threshold $\propto\rho$: sufficiency is a standard consequence of the trace-inverse inequality \cite{Stenberg1995,JuntunenStenberg2009}, while the explicit necessity certificates of Proposition~E appear to be new;
\item the extensions: Clough--Tocher refinements, general smooth domains, steady Navier--Stokes and three dimensions (Theorems~F--I), and the growing-penalty analysis (Theorem~R).
\end{itemize}
That the pressure term is needed for consistency is not new \cite{BurmanHansboLarson2019Stokes}; the mechanism of the leak is the classical consistency error of boundary penalty methods \cite{Babuska1973}, here acting only on the pressure and with vanishing penalty parameter $h/\mu$, which is why the method converges at all.

\subsection*{Scope and related variants}
We analyse the symmetric Nitsche form of \cite{GjerdeScott2024}. The leak comes from the missing pressure term, not from the symmetry of the viscous terms, so we expect it in nonsymmetric variants as well; we have not analysed these, nor penalty-free variants \cite{FreundStenberg1995,Burman2012,BoiveauBurman2016}. Nitsche's method with Robin-type data is treated in \cite{JuntunenStenberg2009}. Multiplier and boundary-correction methods \cite{LiuNeilanOtus2023,BurmanHansboLarson2024,BrambleDupontThomee1972,DupontGuzmanScott2025} reach higher order than $\hG^{3/2}$, at the price of extra unknowns or Taylor corrections.
A common alternative restricts the pressure to $L^2_0$ and adds the plain term $\ip{p_h}{v\cdot n}$, testing the continuity equation only with mean-zero pressures. Because $\div\Vh=\Pih$, this gives only $\div u_h\equiv c$ with $c|\Oh|=\int_{\partial\Oh}u_h\cdot n$ (Proposition~\ref{prop:singular}; compare \cite[(5.9)]{FrachonNilssonZahedi2024}). The constant vanishes exactly when the naive system is solvable, equivalently when the $\CNS$ pressure has zero mean on $\Gh$. With weak imposition nothing forces the flux through $\Gh$ to vanish, even for flux-corrected data, and in our computations $c\ne0$.
Our proofs use Scott--Vogelius elements with $k\ge4$, and $k\ge2$ on Clough--Tocher splits, in two dimensions (with a three-dimensional extension, Theorem~I). Other divergence-free families, and cut or unfitted meshes, are not treated. Gjerde and Scott also treat slip conditions by Nitsche's method \cite{GjerdeScott2022}; we have not examined that method.
%% VERIFY: GjerdeScott2022 full text not read; whether its slip-Nitsche form contains a pressure term is unknown (notes/v5/lit_report.md).

\subsection*{Extensions}
\begin{theoremF}[Theorem~\ref{ct:thm:main}, Proposition~\ref{ct:prop:penalty}]
On the Clough--Tocher refinement of any shape-regular mesh whose boundary edges have comparable lengths, Theorems~A--D, Corollary~B$'$ and the sufficiency part of Proposition~E hold for every $k\ge2$, with no angle or singular-vertex condition. For example, $\CNS$ satisfies $\norm{\ut-u_h}_{H^1}\le C(\hG^{3/2}+h^2)$ for $k=2$ if $\hG\ge h^4$. The penalty necessity is recertified exactly on split stars.
\end{theoremF}

\begin{theoremG}[Section~\ref{gd:sec}]
In two dimensions, Theorems~A--D and Corollary~B$'$ extend to a bounded connected domain whose curved boundary consists of finitely many closed smooth curves, convex or not, with a nonempty polygonal component $\Sigma$ on which the data are imposed strongly. The geometric lemmas use only $C^3$ curves; the results assume the regularity (R1) of the solution, which for generic $f$ requires $\Gamma\in C^{k+1,\alpha}$, and the results that use the leak test field assume $\Gamma\in C^{k+3}$. With several curved components, the correct pressure constant is the \emph{global} boundary mean. $\CNS$ needs one global correction, and a per-component correction is inconsistent.
\end{theoremG}

\begin{theoremH}[Section~\ref{ns:sec}]
For steady Navier--Stokes, assume the discrete linearisation at the exact solution is uniformly stable; we prove this under a small-data condition. Then $\GS$ converges at $H^1$ rate $h/\mu$ with the same leak (slip and energy constants $1$), and $\CNS$ attains $\hG^{3/2}+\hO^k$. No inflow/outflow correction of the convective term is needed for these statements.
\end{theoremH}

\begin{theoremI}[Theorem~\ref{td:main}]
In three dimensions, with an inscribed triangulated polyhedron, Theorems~A--D and Corollary~B$'$ hold with the same rates, assuming a uniform discrete inf-sup constant. That assumption is known for Alfeld (barycentric) splits with $k\ge3$, by the published split-mesh theory combined with a uniform Bogovski\u{\i} bound that we prove; we have not re-audited the published proofs for hidden domain dependence. Only a spherical obstacle is treated. Penalty necessity, sharpness and unsplit meshes remain open in 3D.
\end{theoremI}

\subsection*{Computations}
Section~\ref{sec:numerics} reports computations with up to $N=256$ boundary chords, mostly for $k=4$, with $k=5,6$ and Clough--Tocher $k=2,3$ in Section~\ref{sec:numv3}. They show the following.
\begin{itemize}
\item The $\CNS$ $H^1$ rate decreases monotonically to $3/2$, independently of the pressure, for $k=4,5,6$.
\item The printed method has rate $1$ whenever $G>0$.
\item On a pure-pressure test that isolates the traction response exactly, the slip and energy constants of Corollary~B$'$ approach $1$ for every element tested. For $k=4$ the $H^1$ constants approach the $K$ of Theorem~B$''$; that they do so also for $k=5,6$ and on Clough--Tocher refinements is observed, not proved.
\item In the penalty study, the coercivity threshold $\mu^\ast$ is linear in $h/\hG$ over more than a decade, with $\gamma^\ast=\mu^\ast\hG/h\approx18$--$21$ for $k=4$. It is larger for higher $k$ and on Clough--Tocher refinements.
\item Strongly imposed no-slip shows the $h^{1/2}$ locking on a mesh with two-triangle wall vertices but not on our production meshes, while $\GS$ and $\CNS$ are insensitive to the change of mesh (Section~\ref{num:strong}).
\item The pressure errors follow Theorem~P: rate $\to\frac32$ for $\CNS$, and rate $\to\frac12$ for $\GS$ when $G>0$, carried by the first layer (Section~\ref{pd:sec:num}).
\item A growing penalty behaves as Theorem~R predicts, and flux-corrected data remove the flux floor exactly (Section~\ref{rs:sec:num}).
\end{itemize}
All code, raw outputs and logs are in the accompanying repository (Appendix~\ref{app:repro}).

\subsection*{Organisation}
Section~\ref{sec:setting} fixes the setting. Sections~\ref{sec:geometry}--\ref{sec:erroreq} collect geometric and discrete tools and the error equation. Section~\ref{sec:printed} treats the printed method and identifies its $H^1$ constant. Section~\ref{sec:consistent} treats the pressure-consistent method, proves the lower bound, and treats flux-corrected data and a growing penalty. Section~\ref{pd:sec} treats the pressure, the boundary stress and drag. Section~\ref{sec:penalty} treats the penalty threshold. Sections~\ref{ct:sec}--\ref{td:sec} contain the extensions: Clough--Tocher refinements, general domains, Navier--Stokes and three dimensions. Section~\ref{sec:numerics} reports the computations.
